\documentclass[11pt]{article}\usepackage[margin=1in]{geometry}\usepackage{amsmath,amsthm}\newtheorem{theorem}{Theorem}\newtheorem{remark}{Remark}
\title{Sector-Aware Multichannel Partition Budgets}\author{}\date{21 July 2026}\begin{document}\maketitle
\begin{abstract}We give the exact union and free-energy allocation laws for center, low-dispersion, high-dispersion, cut/orbit, and topology numerators.\end{abstract}
Let $B_{V,c}$ be disjoint or merely covering channel numerators and suppose
\[B_{V,c}(s)\le q_c Z_{G,\mathrm{ref}}(s)V(s).
\]
\begin{theorem}[Typed union law]The union numerator satisfies
\[B_V(s)\le\left(\sum_cq_c\right)Z_{G,\mathrm{ref}}(s)V(s).
\]
For a target $q$ and positive shares $a_c$ summing to one, the allocation $q_c=a_cq$ is exact. With comparison loss $\ell_c$, channel $c$ requires
\[F_c>\ell_c-\log(a_cq).
\]\end{theorem}
\begin{proof}Subadditivity of integrals gives the first statement. Substitution gives the allocation and free-energy formula.\end{proof}
For five equal channels, each free-energy requirement is larger than the undivided target by $\log5$. This prevents five individually plausible bounds from silently overspending one total budget. Ownership labels are retained through the next-state map.
\begin{remark}The allocation theorem is exact. None of the five complete physical channel bounds is yet uniform in volume and retained state.\end{remark}\end{document}
